\section{相关工作}\label{sec:related}

\textbf{可认证的可微层。} SDPRLayers \cite{holmes2025sdpr} 是最接近的工作：它把齐次 QCQP 提升为 Shor SDP，在松弛紧时读出全局解，并用一个能容忍冗余等式约束的隐式选择论证对\emph{原始} KKT 系统求微分。有三点差异。(i)~\emph{约束类}：其理论针对二次\emph{等式}；不等式需要松弛变量，会膨胀 PSD 块并破坏严格互补性，而我们的约束是原生的。(ii)~\emph{SDP 所处位置}：他们在前向传播中对每个样本、每一步求解 SDP；我们前向求解固定维数多项式系统，只在后验证书中异步使用 SDP。(iii)~\emph{全局性依赖什么}：他们依赖前向松弛；我们依赖一个独立的 $14\times14$ 松弛，因此紧性失效只会使我们失去证书，而不会失去前向求解。

\textbf{代数全局四元数求解器。} QPEP \cite{wu2022qpep} 把 PnP、PnL、手眼标定、点到面配准、圆锥位姿、基于距离的相对位姿与 Stewart 平台运动学统一为“二次位姿估计问题”，消去 $\mathbf t$，并用 Gr\"obner 基枚举所有 KKT 根（至多 40 个候选）。我们的工作为 QPEP 增加不等式与梯度：保留其消元与枚举，并加入仿射约束、ALM 循环、FB 微分与半定后验证书。在 QPEP 的意义上，“全局最优”指给定完全枚举时的最佳驻点，这不是凸证书。

\textbf{可微凸优化层。} OptNet \cite{amos2017optnet} 通过内点求解的 KKT 系统对 QP 求微分；CVXPYLayers \cite{agrawal2019cvxpylayers} 对锥规划做同样的事；dQP \cite{magoon2025dqp} 将求解器与导数解耦，并表明活跃集决定二者。这些方法假设凸性、已辨识的活跃集以及 LICQ 或严格互补性；我们的可行集位于 $\mathbb{S}^3$ 上，LICQ 在四接触与三个面外自由度处失效，严格互补性在擦边接触处失效。我们借用 dQP 的思想——把导数约简为由活跃集确定的小系统——并加入秩鲁棒正则化与极限理论。

\textbf{局部位姿层与非光滑微分。} BPnP \cite{chen2020bpn} 通过隐式微分穿过 PnP 求解反向传播；EPro-PnP \cite{chen2022epropnp} 用 Monte Carlo 积分位姿后验；Theseus \cite{pineda2022theseus} 在流形上提供带隐式或展开导数的可微非线性最小二乘。三者都是局部的，支持惩罚但不支持硬约束。Fischer--Burmeister 函数 \cite{fischer1992} 及其光滑变体 \cite{kanzow1996} 把互补性转化为半光滑方程 \cite{facchinei2003}；非凸 ALM 是经典方法 \cite{bertsekas1982,rockafellar1998}，精确惩罚变体见 \cite{fletcher1987}；\cite{bolte2021nonsmooth} 则是我们陈述命题~\ref{prop:limit} 极限的框架。表~\ref{tab:positioning} 在附录中给出比较。
